\documentclass[12pt]{article}
\usepackage[applemac]{inputenc}

\usepackage[makeroom]{cancel}
\usepackage{amsmath}
\usepackage{amssymb}

%\textwidth=16.7truecm \textheight=22.8truecm \topmargin=-1.truecm
%\oddsidemargin=-1.truecm \evensidemargin=-1.truecm

\begin{document}

\pagestyle{empty}

\centerline{\Large \bf NORDAN 2016: Several Complex Variables}
\vspace{3mm}

\centerline{\large \bf Location: room ??}

\centerline{\large \bf Abstracts}

%
% Daniel Barlet <daniel.barlet@iecn.u-nancy.fr>
% Moulay Taib Belghiti <belghititaib@yahoo.fr>
% Judith Brinkschulte <brinkschulte@math.uni-leipzig.de>
% (preferebly on Thursday or Friday).
% Miroslav Englis <englis@math.cas.cz>
% Mattias Jonsson <mattiasj@umich.edu>
% Frank Kutzschebauch <frank.kutzschebauch@math.unibe.ch> 
% J�n Ing�lfur Magn�sson <jim@hi.is>
% Francine Meylan <francine.meylan@unifr.ch>
% (preferably on Saturday)
% David Witt Nystr�m <danspolitik@gmail.com>
% (only on Friday)
% Dan Popovici <popovici@math.univ-toulouse.fr>
% Alexander Rashkovskii <alexander.rashkovskii@uis.no> 
% Martin Sera <sera@math.uni-wuppertal.de>
% Margaret Stawiska-Friedland <stawiska@umich.edu>
% Maria Trybula <maria.h.trybula@gmail.com>
% 

%\end{document}




\begin{center}\large
{A note on  some fiber-integrals  $\int_{f = s}\
  \rho.(\omega/df)\wedge(\overline{\omega/df}) $} 
\\ \normalsize \vspace{0.2cm}
{\bf  Daniel Barlet }\\{Universit\'e de Lorraine
and Institut Universitaire de France}
\end{center}


\noindent 
We remark that the study of a fiber-integral of the type 
  $$F(s) := \int_{f = s}\ \rho.(\omega/df)\wedge(\overline{\omega/df}) $$
   either in the local case where $\rho \equiv 1$ around $0$ is
   ${\mathcal C}^{\infty}$ and compactly supported near the origin
   which is a singular point of $\{f = 0\}$ in ${\mathbb C}^{n+1}$, or in a
   global setting where $ f : X \to D$ is a proper holomorphic
   function on a complex manifold $X$, smooth outside $\{f = 0\}$ with
   $\rho \equiv 1$ near $\{f = 0\}$, for given holomorphic
   $(n+1)-$forms $\omega$ and $\omega'$, that a better control on the
   asymptotic expansion of $F$ when $s \mapsto 0$, is obtained by
   using the Bernstein polynomial of the ``frescos'' associated to $f$
   and $\omega$ and to $f$ and $\omega'$ (a fresco is  a ``small''
   Brieskorn module corresponding to the differential equation deduced
   from the Gauss-Manin system of $f$ at $0$) than to use the
   Bernstein polynomial of the full Gauss-Manin system of $f$ at the
   origin. We illustrate this in the local case in some rather simple
   (non quasi-homogeneous) polynomials, where the Bernstein polynomial
   of such a fresco is explicitly evaluated. 





%%% macros for Belghiti
\def \nat{\mathbb{N}}
\def \com{\mathbb{C}}
\def \ree{\mathbb{R}}
\def \ppoe{\leqslant}
\newcommand{\parentheses}[1]{\ensuremath{\left ( #1 \right )}}
\newcommand{\modul}[1]{\ensuremath{ \left \vert #1\right \vert}}
\newcommand{\norm}[1]{\ensuremath{ \left \Vert #1\right \Vert}}
\newcommand{\braces}[1]{\ensuremath{\left \lbrace #1 \right \rbrace}}


\begin{center}\large
{Degrees of polynomial appoximation in \\ 
holomorphic Carleman classes} 
\\ \normalsize \vspace{0.2cm}
{\bf Moulay Taib Belghiti }\\{ Ibn Tofa{\"\i}l University, Morocco}
\end{center}
\noindent 
In this talk, we extend results of M.S. Baouendi and C. Goulaouic
(Ann. Inst. Fourier, $1971$ ; Trans. Amer. Math. Soc, $1974$),
obtained for compacts of $\ree^{N}$ with analytic boundary.
If $K$ is a compact of $\com^{N}\backsimeq\ree^{2N}$, ($N\geq 1$),
$\mathcal{H}_{M}(K)$ is the space of $\overline{\partial}$-Whitney
jets on $K$ which are of class $\{M\}$, where $M(t)=t^{t}e^{t\mu(t)}$,
$t>>0$ and $\mu$ belongs to a Hardy field.
We prove that a jet $F:=\parentheses{F^{\alpha}}_{\alpha\in\nat^{2N}}
\in\mathcal{H}_{M}(K)$ if and only if there exists a constant $C>0$,
such that 
\begin{equation*}\label{equation 1}
\displaystyle\lim_{n\rightarrow\infty}d_{n}(F^{\alpha},K)\exp({C \overline{\omega}_{K,{M}}}(n))=0,\quad \text{ for all}\quad \alpha\in\nat^{2N},
\end{equation*}
where $d_{n}(\cdot,K)$ is the distance, for the uniform norm on $K$  to the complex vectorial space of polynomials of degree at most $n$, and where $\overline{\omega}_{K,{M}}$ is a weight depending on the class $\{M\}$ and $K$.
If $K$ is Whitney-regular 
\begin{multline*}
\mathcal{H}_{M}(K)\backsimeq \big
\{f\in\mathcal{E}^{\infty}(K)\cap\mathcal{O}(\dot{K}): \\
\exists  C>0,\;\exists\rho>0,\; \norm{{D^{\alpha}}f}_{K}\ppoe
  C\rho^{\modul{\alpha}}M(\modul{\alpha}),\;(\forall
  \alpha\in\nat^{N})\big\}. 
\end{multline*}
 In this situation, $f\in\mathcal{H}_{M}(K)$ if and only if
 $\displaystyle\lim_{n\rightarrow\infty}d_{n}(f,K)e^{C
   \overline{\omega}(n)}=0$, where $C>0$  and $\overline{\omega}$ is a
 weight depending on \{M\}. 
Finally, we annonce similar results in the situation where $K$ is a
compact of some Stein manifold. A crucial role is played by a new
geometric criteria: the {\L}ojasiewicz-Siciak condition for the Green
function of $K$. 
This is a joint work with Boutayeb El Ammari and Laurent P.~Gendre.


\vspace{0.7cm}




\begin{center}\large
 On the nonvanishing of abstract \\ Cauchy-Riemann cohomology groups
\\ \normalsize \vspace{0.2cm}
{\bf  	Judith Brinkschulte}\\{University of Leipzig, Germany}
\end{center}
\noindent 
In this talk, I will discuss a joint work with C.D.\ Hill and
M.\ Nacinovich. Namely we considered abstract smooth CR structures on a
smooth
manifold M, and the associated global abstract cohomology groups
$H^{p,q}(M)$, which are the analogues of the Dolbeault cohomology
groups.
We have shown that some of these global cohomology groups must be
infinite
dimensional, or non Hausdorff, whenever one has a certain condition on
the
Levi form of the CR structure. What makes our results curious is that
the
required condition on the Levi form needs to be satisfied only at a
single
(mircro-local) point on M; yet the conclusion is global.

\vspace{0.7cm}

\newpage

\begin{center}\large
{High-power asymptotics of weighted harmonic Bergman kernels}
\\ \normalsize \vspace{0.2cm}
{\bf   Miroslav Engli{\v s} }\\{
Institute of Mathematics
of the Czech Academy of Sciences and \\
Silesian University in Opava,  Czech Republic
}
\end{center}

\noindent
The asymptotics of the weighted Bergman kernels with respect to
the weight $|r|^\alpha$, where $r$ is a defining function for a
smoothly bounded strictly pseudoconvex domain and $\alpha\to+\infty$,
play prominent role in mathematical physics (Berezin quantization)
as well as in complex geometry (Donaldson's balanced metrics);
the standard tool for their derivation is the famous description
of the boundary singularity of the Bergman kernel due to Fefferman,
combined with a construction due to Forelli and Rudin.
The talk will describe why it is noteworthy to study the analogous
asymptotics also for the Bergman kernels for harmonic functions,
and will give a complete answer for the case of radial weights
on the ball and horizontal weights on the upper half-space.
The proofs actually proceed by relating the problem to the
holomorphic case mentioned above, but on a different domain.

\vspace{0.7cm}



\begin{center}\large
{Invariant measures for birational surface maps\\  defined over number fields} 
\\ \normalsize \vspace{0.2cm}
{\bf  Mattias Jonsson }\\{University of Michigan and \\
Chalmers University of Technology and the University of Gothenburg}
\end{center}

\noindent
Iterating a birational selfmap of the complex projective plane can
lead to very interesting dynamics. When the dynamical degree is larger
than one, it is expected that there exists a unique measure of maximal
entropy, and that this measure is obtained as the intersection of two
positive closed currents. Unfortunately, this intersection is only
known to be well-defined under hypotheses that are often hard to
verify. I will report on joint work with Paul Reschke, where we prove
that if the map has rational (or algebraic) coefficients, then the
intersection is well defined and the complex dynamics is under
control. 


\vspace{0.7cm}


\newpage 

\begin{center}\large
An Oka principle for simultaneous \\ standardization of n-tuples of points
\\ \normalsize \vspace{0.2cm}
{\bf Frank Kutzschebauch}\\{University of Bern, Switzerland}
\end{center}
\noindent 
It is an easy exercise to show that the holomorphic automorphism group
$Aut_{hol} ({\mathbb C}^n)$ $n\ge 2$ acts transitively on  ordered N-tuples
$(z_1, z_2, \ldots  z_N)$ of pairwise disjoint points in ${\mathbb C}^n$. If
the  points depend holomorphically on a Stein parameter $w \in W$ we
ask whether the automorphism of ${\mathbb C}^n$ moving the N-tuple $(z_1(w),
z_2(w), \ldots  z_N(w))$ to a fixed N-tuple $(z_1, z_2, \ldots  z_N)$
can be chosen holomorphically depending on the parameter $w$. We prove
an Oka principle saying that the obstruction for this is of purely
topological nature. Our theorem (which is true not only for ${\mathbb C}^n$ but
for any Stein manifold with the density property) can be interpreted
as a result similar  to Grauerts Oka principle, but instead for  the
principal bundle of a complex Lie group in Grauerts case for  a
principal bundle of  certain infinite-dimensional Frechet groups in
our case. 
This is a joint work with Ramos Peon.



\begin{center}\large
Intersection theory and relative fundamental classes
\\ \normalsize \vspace{0.2cm}
{\bf J\'on Ing\'olfur Magn\'usson}\\{University of Iceland}
\end{center}


\noindent 
We will describe an intersection theory for analytic cycles in a
complex manifold. Main properties will be explained, in particular
with respect to analytic families of cycles and their relative
fundamental classes. 
This is a joint work with Daniel Barlet.


\vspace{0.7cm}




%% Macros for Meylan

\def\dbl{[\hskip -1pt[}
\def\dbr{]\hskip -1pt]}
\newcommand{ \al}{\alpha}
\newcommand{\ab}[1]{\vert z\vert^{#1}}

\begin{center}\large
Automorphism groups of Levi degenerate hypersurfaces in $\mathbb C^3$
\\ \normalsize \vspace{0.2cm}
{\bf Francine Meylan}\\{University of Fribourg, Switzerland}
\end{center}

\noindent 
We give a classification of  smooth real hypersurfaces of finite
Catlin multitype in $\mathbb C^3$ which admit  
nonlinear 
infinitesimal automorphisms.
The results are complete on the level of weighted homogenous
polynomial models. 
As a consequence,  we prove a sharp 1-jet determination result in the
general case. 
We  also  identify a common source of such vector fields.
They all arise by pulling back a symmetry of a suitable 
hyperquadric in $\mathbb C^K$, $K \geq 3,$ by a holomorphic mapping
from  $\mathbb C^3$  to $\mathbb C^K.$ 
This is a joint work with Martin Kolar.

\vspace{0.7cm}



\begin{center}\large
{Duality between the pseudoeffective and \\ the movable cone on a
  projective manifold.}  
\\ \normalsize \vspace{0.2cm}
{\bf   	David Witt Nystr\"om}\\
{Chalmers University of Technology and the University of Gothenburg}
\end{center}



\noindent 
The structure of projective algebraic manifolds is to a large extent
governed by the geometry of its cones of divisors or curves. In the
case of divisors, two cones are of primary importance: the cone of
ample divisors and the cone of effective divisors (and the closure of
these cones as well). These cones have natural transcendental
analogues, namely the cone of K�hler classes (called the K�hler cone)
and the cone of pseudoeffective $(1,1)$-classes (called the
pseudoeffective cone).   
 
I will discuss my recent proof of a conjecture of
Boucksom-Demailly-Paun-Peternell which says that on a projective
manifold the pseudoeffective cone is dual to the cone of movable
classes. A celebrated result of Demailly-Paun showed that the K�hler
cone is determined by analytic cycles and the Hodge structure of the
manifold. A consequence of the duality theorem is that on a projective
manifold the pseudoeffective cone is similarly determined by the
analytic cycles and the Hodge structure of the manifold together with
all its modifications. 



\vspace{0.7cm}





\begin{center}\large
Positivity and Duality on Compact Complex Manifolds
\\ \normalsize \vspace{0.2cm}
{\bf Dan Popovici }\\{Paul Sabatier University, Toulouse, France}
\end{center}


\noindent 
Given a compact complex manifold $X$ of dimension $n$, we shall
explain how the duality between Demailly's pseudoeffective cone of
Bott-Chern cohomology classes of closed positive $(1,1)$-currents and our
Gauduchon cone of Aeppli cohomology classes of $(n-1)^{st}$ powers of
Gauduchon metrics can be used, together with a natural way of estimating
from below the integral of a quantity involving the solution of a
Monge-Amp\`ere equation provided by Yau's theorem, to completely prove the
qualitative part and give a partial affirmative answer to the quantitative
part of Demailly's Transcendental Morse Inequalities Conjecture for
differences of two nef classes. If time permits, we shall explain how a
complete solution to this conjecture would contribute to tackling the
conjecture we have proposed in the non-K\"ahler context predicting that
every $\partial\bar\partial$-manifold should carry a balanced metric.


\vspace{0.7cm}



\begin{center}\large
{Analytic and plurisubharmonic singularities} \\
{\bf  Alexander Rashkovski{\v\i}  }\\{University of Stavanger, 
Norway}
\end{center}

\noindent
The subject of the talk is how good one can approach isolated
plurisubharmonic singularities by analytic ones. 
The following two (related) open problems will be discussed: existence
of a plurisubharmonic function with the zero Lelong number and
positive residual Monge-Amp�re mass, and control over the residual
Monge-Amp�re masses of analytic approximations. 


\vspace{0.7cm}




\begin{center}\large
{On vanishing results for $L^2$-Dolbeault cohomology on pseudoconvex
complex spaces }
\\ \normalsize \vspace{0.2cm}
{\bf  Martin Sera}\\{University of Wuppertal, Germany}
\end{center}


\noindent 
Using Grauert's bumping method, it was proven that $L^2$-Dolbeault
cohomo\-logy groups of a strictly pseudoconvex bounded domain in a
complex manifold are isomorphic to Dolbeault cohomology groups with
respect to locally square integrable forms. We generalize this to
strictly pseudoconvex domains in complex spaces. For this, we prove
vanishing theorems for $L^2$-Dolbeault cohomology groups on
pseudoconvex domains of analytic sets in $\mathbb{C}^n$, which
generalize results by T. Ohsawa, W. Pardon and M. Stern. 


\vspace{0.7cm}




\begin{center}\large
A characterization of polynomials in complex and non-archimedean dynamics
\\ \normalsize \vspace{0.2cm}
{\bf Margaret Stawiska-Friedland}\\{Mathematical Reviews, Ann Arbor, USA}
\end{center}

\noindent 
In 1960s Hans Brolin initiated systematic application of
potential-theoretic methods in the dynamics of holomorphic
polynomials. Among other things, he proved the now-famous
equidistribution theorem: for a polynomial $f$ of degree greater than
$1$ the preimages, under successive iterates of $f$, of a Dirac
measure at an arbitrary point of the complex plane (except at most two
so-called exceptional points) converge weakly to the equilibrium
measure of the Julia set for $f$. In 1980s a similar result (about
convergence of preimages of quite general probabilistic measures) was
proved for a rational map $f$ of degree greater than $1$. The limit
measure obtained in this case (called the balanced measure) is also
supported on the Julia set for $f$, but does not have to be its
equilibrium measure. In fact, A.O. Lopes proved (using dynamical
properties of Julia sets) that equality of these two measures (under
suitable assumptions on $f$, also making precise the notion of the
equilibrium measure for the Julia set) implies that $f$ is a
polynomial. In 2010 we obtained  a proof of Lopes's theorem (under
slightly weaker assumptions) using only  classical and weighted
potential theory. In this talk I will present a recent extension of
this result, namely a characterization of polynomials 
among rational functions,
up to rational functions having potentially good reductions
as exceptions, on the projective line
over an algebraically closed field of any characteristic
that is complete with respect to a non-trivial and possibly non-archimedean
absolute value. I will introduce basic notions of non-archimedean
dynamics and discuss possible cases. 
This is joint work with Yusuke Okuyama from Kyoto Institute of Technology.

\vspace{0.7cm}


\begin{center}\large
{Boundary behavior of invariant functions on planar domains} 
\\ \normalsize \vspace{0.2cm}
{\bf  Maria Trybu{\l}a}\\{Adam Mickiewicz University, 
Pozna\'n, Poland}
\end{center}

\noindent
In the talk I will show the precise  behavior of Carath\'eodory,
Kobayashi and Bergman metrics and the corresponding
distances near smooth boundary
points of planar domains under different assumptions of
regularity. Some applications will be given. 
The lecture is based on a joint paper with N. Nikolov and L. Andreev

\end{document}

